%% ============================================================
%%  Section: Identification Strategy
%%  Depositor Discipline in Modern Banking (2026-04-07 draft)
%%
%%  Changes from 20260404 version:
%%    - Removed all [@@@ ...] placeholders and addressed each one.
%%    - Toned down "silver bullet" / overselling language throughout.
%%      "instrument" → "natural experiment"; softened "precisely identified"
%%      and "orthogonal" claims.
%%    - Removed all \paragraph{} headers; content flows as prose.
%%    - Fewer em-dashes; replaced with commas or parentheses.
%%    - Staggered-adoption paragraph shortened; removed Callaway/Sun/
%%      de Chaisemartin citations and removed stacked-cohort robustness
%%      claims that are not in the code.
%%    - Phase-in election paragraph kept (mechanism-test role) but shortened.
%%    - Threats to Identification: collapsed to one paragraph; removed
%%      Anticipation paragraph and General Equilibrium paragraph per user
%%      instruction.
%%    - Signal Validation: shortened by ~half; removed \paragraph{} headers
%%      and detailed coefficient recitation; retained key findings.
%%    - Removed the BIB ENTRIES comment block at end of file.
%%
%%  Cross-references:
%%    eq:dayone  (sec:background)
%%    tab:descriptive_sumstats_cecl, tab:cecl_predictors_cross_section
%%    fig:cecl_equity_distribution
%%    fig:es_large_banks_market_cds
%%    sec:results, sec:results:heterogeneity
%% ============================================================

\section{Identification Strategy}
\label{sec:identification}

Our goal is to identify the causal effect of a credit-risk disclosure on
depositor behavior.  The empirical obstacle is familiar: riskier banks
carry fewer uninsured deposits in equilibrium, so cross-sectional comparisons
conflate depositors' information responses with pre-existing sorting and
equilibrium co-determination \citep{covitz2004market, chen2022bank}.
We address this problem by exploiting the CECL Day-One adjustment as a
natural experiment that delivers a well-defined information shock, one whose
timing is set by regulation and whose magnitude reflects the bank's assessment
of lifetime expected losses on its existing loan portfolio---incorporating both
historical lending decisions and current economic conditions at adoption.  The
cross-sectional intensity of the shock is therefore largely fixed by past
underwriting rather than by the deposit dynamics we study, a claim we treat as
testable rather than definitional and probe with extended pre-trends,
residualized treatments, and the within-bank composition design below.

\subsection{Empirical Design}
\label{sec:identification:did}

Our baseline specification is a difference-in-differences estimator that
exploits cross-sectional variation in CECL exposure and the common
regulatory adoption date:
\begin{equation}
  y_{it}
  = \alpha_i + \gamma_t
  + \beta \cdot \mathrm{Post}_{it} \times \mathrm{CECL}_i
  + \mathbf{X}_{it-1}'\boldsymbol{\delta}
  + \varepsilon_{it},
  \label{eq:did}
\end{equation}
where $y_{it}$ is the outcome for bank $i$ in quarter $t$; $\alpha_i$ and
$\gamma_t$ are bank and calendar-quarter fixed effects; $\mathrm{Post}_{it}$
equals one for all quarters at or after bank $i$'s CECL adoption date;
$\mathrm{CECL}_i$ is either the continuous \texttt{CECL Adj} measure or the
binary \texttt{High CECL} indicator, both defined in
Section~\ref{sec:background}; and $\mathbf{X}_{it-1}$ is a vector of lagged
controls (log assets, equity-to-assets ratio, interest expense-to-assets).
Bank fixed effects absorb all time-invariant differences across institutions,
including balance-sheet composition, business model, geographic market, and
charter type, so that $\hat{\beta}$ is identified from within-bank changes
in outcomes over time.  Quarter fixed effects absorb any aggregate
macroeconomic force common to all banks in a given quarter, including the
Federal Reserve's tightening cycle, aggregate deposit normalization, and any
nationwide reaction to the 2023 banking stress.

To evaluate the parallel-trends assumption and trace the dynamics of depositor
responses, we also estimate an event-study specification:
\begin{equation}
  y_{it}
  = \alpha_i + \gamma_t
  + \sum_{\substack{k=-10 \\ k \neq -1}}^{10}
    \beta_k \cdot \mathbf{1}\!\left\{\tau_{it} = k\right\}
    \times \mathrm{CECL}_i
  + \mathbf{X}_{it-1}'\boldsymbol{\delta}
  + \varepsilon_{it},
  \label{eq:eventstudy}
\end{equation}
where $\tau_{it} = t - t_i^{*}$ is event time in quarters relative to bank
$i$'s CECL adoption date, and $k = -1$ is the omitted reference period.
Coefficients $\{\hat{\beta}_k\}_{k=-10}^{-2}$ test for pre-adoption
differential trends between high- and low-CECL banks, with $k=-1$ normalized
to zero as the reference period; coefficients $\{\hat{\beta}_k\}_{k \geq 0}$
trace the cumulative post-adoption effect.  A pattern of flat pre-trends
followed by a post-adoption break supports the interpretation that observed
differences are caused by the CECL disclosure rather than pre-existing
divergence.  Our baseline window spans eleven quarters on each side of
adoption, so the pre-period begins in 2020Q3 and excludes the COVID-19
disruption quarters (2020Q1--2020Q2), whose aggregate liquidity dislocations
complicate the interpretation of individual pre-trend coefficients.  Because
this choice could in principle conceal unfavorable history, we
also re-estimate the event study on the full panel back to 2018Q1, both
including and excluding the COVID quarters, and report joint tests of the
pre-adoption coefficients (Section~\ref{sec:robustness:pretrends}).


\subsection{Treatment Exogeneity}
\label{sec:identification:predetermination}

For $\hat{\beta}$ in \eqref{eq:did} to have a causal interpretation,
the CECL treatment must be uncorrelated with potential outcomes absent the
disclosure.  This requires that neither the timing of the shock nor its
magnitude was determined by, or correlated with, the deposit dynamics we
study.

The 2023Q1 adoption date was set by FASB's ASU 2016-13 and enforced by the
federal banking agencies.  Institutions subject to the community-bank rollout
had no ability to accelerate or delay adoption to coincide with favorable
deposit market conditions, so the adoption date is unrelated to any
bank-specific information about contemporaneous deposit flows.

The size of the Day-One adjustment reflects the cumulative composition of the
loan portfolio as it stood on the adoption date, itself the product of
lending decisions made years or decades prior.  The CECL charge is large when
a bank's portfolio contains loans with long remaining maturities, weak
borrower credit quality, or thin collateral coverage, characteristics fixed
by prior underwriting decisions rather than 2023 deposit conditions.  The
cross-sectional regression in Table~\ref{tab:cecl_predictors_cross_section}
confirms that the CECL adjustment is predicted by NPL ratios, size, and
capital in the expected directions, and---in an expanded specification that
adds deposit composition, eight-quarter profitability and asset-quality
trends, and local deposit-market concentration---by the time-deposit and
brokered shares of funding, with a maximum adjusted $R^2$ of roughly 1.4
percent.  We are careful about what this low explanatory power does and does
not establish.  It does not prove exogeneity: a treatment could be orthogonal
to every observable and still correlate with unobserved determinants of
deposit flows.  What it establishes is narrower and still useful: high-CECL
banks are not simply the observably weak, deposit-fragile, or stress-exposed
banks relabeled, so the treatment delivers variation that is largely
unspanned by the characteristics a depositor---or an econometrician---could
have used to sort banks in 2022.  We therefore lean on three complementary
pieces of evidence rather than on the $R^2$ itself: flat extended
pre-trends, robustness to a treatment \emph{residualized} on the full
expanded predictor set (so that the identifying variation is by construction
orthogonal to the observables), and the within-bank composition design of the
triple-difference, which is immune to bank-level selection on unobservables.

Two anticipation margins deserve separate treatment.  The standard itself was
public from June 2016, so banks and depositors knew adoption was coming;
what was not knowable in advance was the \emph{magnitude} of any given bank's
Day-One charge, which depended on modeling choices, portfolio seasoning, and
macroeconomic forecasts finalized only at adoption.  Anticipation of the
event date without knowledge of treatment intensity does not bias the
difference-in-differences estimate; anticipation of intensity would appear as
pre-adoption divergence, which the extended event studies are designed to
detect.  On the bank side, lenders could in principle have adjusted deposit
pricing ahead of adoption; the lagged interest-expense control absorbs such
pre-positioning, and the pre-period coefficients on funding costs are flat.


\subsection{Information Content, Not Cash Flows}
\label{sec:identification:mechanism}

A central advantage of our design is the separation of information from
concurrent changes in bank financial condition.  Prior depositor-discipline
tests typically exploit realized distress events, rating downgrades, NPL
spikes, or near-failure episodes, in which new information arrives
simultaneously with deterioration in actual cash flows or liquidity
\citep{flannery1996evidence, martinez2001depositors, iyer2016tale}.
In such settings, the response to information is inseparable from the
response to contemporaneous financial deterioration.  The CECL Day-One
adjustment resolves this confound: as discussed in
Section~\ref{sec:background}, it records a pure accounting charge that does
not trigger a cash transfer, alter borrower repayment terms, reduce bank
liquidity, or change the risk of the underlying loans.  A depositor who
withdraws funds or demands a higher rate after observing a large CECL
adjustment is responding to newly disclosed information about credit-risk
composition, not to a concurrent worsening of the bank's actual financial
condition.

What distinguishes the CECL Day-One adjustment from routine Call Report
disclosures---which are also mandatory, quantified, and audited---is the
predetermined adoption date and the forced aggregation of forward-looking
lifetime credit losses into a single disclosed figure.  Unlike NPL ratios or
allowance levels reported each quarter at management's discretion within
accounting rules, the Day-One adjustment requires banks to reveal their full
expected-loss assessment in one number at a calendar date set by regulation,
not by management's disclosure preferences.  This eliminates the endogeneity
problem of voluntary disclosure events and the behavioral-constraint problem
of supervisory rating changes.

\subsection{The 2023 Banking Stress and Identification}
\label{sec:identification:stress}

The principal threat to identification is the coincidence of the primary
adoption quarter (2023Q1) with the failures of Silicon Valley Bank and
Signature Bank in March 2023.  Calendar-quarter fixed effects absorb only the
\emph{common} component of that shock.  The sharper version of the concern is
indirect: if exposure to the 2023 stress varied across banks with
characteristics that correlate, even weakly, with the CECL adjustment, then
the differential post-2023Q1 dynamics of high-CECL banks could reflect
differential stress exposure rather than the disclosure.  Depositors ran in
2023 on banks with large uninsured shares and unrealized securities losses
\citep{caglio2024depositors, jiang2024monetary}; quarter fixed effects do not
rule out high-CECL banks responding to that episode through such balance-sheet
channels.  We confront this concern with three pieces of evidence, developed
in Section~\ref{sec:robustness:stress}, beyond the timing fact established in
Section~\ref{sec:background} that the Day-One figures became public only after
the March panic.

First, ex ante exposure to the stress channels is essentially uncorrelated
with treatment: in Table~\ref{tab:cecl_predictors_cross_section}, unrealized
securities losses and the uninsured deposit share---the two characteristics
most predictive of 2023 run exposure---have no meaningful relationship with
the CECL adjustment.  Second, we relax the common-shock restriction directly,
interacting the full vector of 2022Q4 bank characteristics that governed 2023
stress exposure with calendar-quarter fixed effects, so that each
characteristic carries its own free time path.  The within-bank composition
shift survives this saturation essentially intact.  The \emph{level} of the
uninsured decline attenuates once deposit-composition variables receive their
own time paths---a pattern we discuss at length, because deposit composition
is partly an outcome of the channel under study and controlling for its time
path absorbs part of the treatment effect itself.  Third, we exploit the
staggered rollout: banks adopting in 2022Q3--2022Q4 and 2023Q2--2023Q4
experienced the March 2023 shock at different event times than the primary
cohort.  Event studies in calendar time show responses aligned with each
cohort's own adoption quarter, and the difference-in-differences estimated
without the 2023Q1 cohort carries the same sign, although the small number of
treated banks outside the primary cohort limits power.  We also report a
placebo using not-yet-adopted banks, and we are transparent that it is not
clean: among late adopters, 2023Q1 deposit changes correlate with the size of
the still-undisclosed future adjustment, a pattern consistent with the March
stress loading partly on the same credit-risk fundamentals that CECL later
capitalized.  This is precisely why the characteristics-by-quarter controls
and the bank-by-quarter triple-difference, rather than any single placebo,
carry the identification burden for the stress confound.

A secondary concern is differential rate sensitivity: if high-CECL banks were
more exposed to the Federal Reserve's 2022--2023 tightening cycle, their
deposit dynamics would reflect rate pass-through rather than information.
The lagged interest expense-to-assets control absorbs pre-existing
differences in pricing strategy, calendar-quarter fixed effects absorb
aggregate rate variation, and the characteristics-by-quarter specifications
allow rate exposure proxies (size, deposit mix) their own time paths.

\subsection{Signal Validation: Market Responses at Large-Bank Adoption}
\label{sec:identification:signal}

A necessary condition for our interpretation is that the CECL adjustment
contained information about bank credit quality that was not already reflected
in prices.  We test this using the cohort of large, publicly traded banks that
adopted CECL in 2020Q1, for which equity prices and CDS spreads are observable
around the adoption date.  We do not use this cohort for
our main tests because the 2020Q1 adoption date coincides with the onset of
the COVID-19 pandemic, which overwhelmed any deposit-market signal with
aggregate liquidity demand and emergency policy interventions.

Figure~\ref{fig:es_large_banks_market_cds} reports event-study estimates of
the differential market response around CECL adoption, using the continuous
\texttt{CECL Adj} ratio as treatment intensity with bank and
calendar-month fixed effects.  Panel~A shows equity market capitalization:
the pre-adoption coefficients across nine lead months are flat and
statistically indistinguishable from zero, ruling out systematic anticipation
by equity investors.  At the adoption month, market capitalization drops
sharply and persistently for banks with larger CECL adjustments.  Panel~B
shows CDS spreads: the pre-adoption pattern is similarly flat, and spreads
widen materially at adoption and remain elevated in subsequent months.
Both markets revised assessments sharply at the moment of disclosure, not
before---though we note that banks with larger CECL adjustments (i.e.,
riskier loan portfolios) also faced larger COVID-19 credit shocks in 2020Q1,
an observationally equivalent alternative explanation for the adoption-month
break.  The flat pre-adoption coefficients confirm that equity and CDS markets
were not systematically anticipating the disclosure, but the pandemic's
simultaneous arrival prevents cleanly attributing the adoption-month movement
to CECL alone.  This evidence is consistent with the CECL adjustment conveying
credit-risk information in a setting where market prices are observable; we
view it as suggestive rather than definitive, and note that the
community-bank depositor tests in Section~\ref{sec:results} do not rely on
this cohort for identification.
